\documentclass[11pt]{article}

% ---- Packages ----
\usepackage[margin=1in]{geometry}
\usepackage{amsmath,amssymb}
\usepackage{graphicx}
\usepackage{booktabs}
\usepackage{hyperref}
\usepackage{setspace}
\onehalfspacing

% Point LaTeX at the figures folder so \includegraphics needs no path prefix
\graphicspath{{figures/}}

% ---- Document ----
\begin{document}

\title{Posture and Sustainment Optimization Under Adversarial Uncertainty}
\author{}
\date{}
\maketitle

\tableofcontents
\newpage

% =========================================================
% SECTION: Experiment 1 — Greedy Baseline Characterization
% =========================================================

\section{Experiment 1: Greedy Baseline Characterization}
\label{sec:exp1}

\subsection{Motivation}

Before evaluating scenario-weighted stochastic optimizers, it is necessary
to establish what a myopic, value-maximizing greedy strategy achieves and,
critically, where it fails. A greedy placement policy constitutes the
expected-value benchmark against which the Composite Expected Value (CEV)
optimizer is measured in Experiment~2. This section characterizes greedy
performance across four posture metrics over a 10-step sustainment horizon,
then quantifies its structural vulnerability to non-uniform threat
distributions---the gap whose elimination motivates distributionally robust
optimization.

%------------------------------------------------------------
\subsection{Experimental Setup}
\label{sec:exp1:setup}

The simulation environment consists of $|\mathcal{A}| = 20$ assets distributed
across $|\mathcal{L}| = 5$ named theater locations drawn from
Indo-Pacific basing sites (Kadena AB, Andersen AFB, MCAS Iwakuni,
Camp H.M.\ Smith, and Diego Garcia), each with a per-location capacity
of $c = 5$ and a pre-assigned strategic value
$v_\ell \in [0.78,\, 0.95]$.
Asset types are drawn uniformly from five categories (aircraft, fuel depots,
maintenance crews, munitions, and medical assets), with initial readiness
rates sampled from $\mathcal{U}(0.4, 1.0)$ and maintenance windows
from $\mathcal{U}(1, 90)$ days.

A fixed per-step degradation rate $\delta = 0.08$ is applied to every
asset at each of $T = 10$ discrete time steps, replacing the stochastic
degradation draws used in general simulation to isolate the effect of the
placement and sustainment policy from noise. All results are averaged over
$N_\mathrm{seed} = 10$ independent random seeds, with standard deviations
reported throughout.

%------------------------------------------------------------
\subsection{Greedy Placement Algorithm}
\label{sec:exp1:algorithm}

The greedy placement policy $\pi_G$ assigns assets sequentially to the
highest-value location with remaining capacity:
\begin{equation}
    \pi_G(a) = \operatorname*{arg\,max}_{\ell \in \mathcal{L}}
    \; v_\ell
    \quad \text{subject to } |\{a' : \pi_G(a') = \ell\}| < c.
    \label{eq:greedy}
\end{equation}
Assets are processed in the order they appear in the initial state. Because
all five locations share the same capacity $c = 5$, and 20 assets exactly
fill four locations, $\pi_G$ deterministically assigns assets to the four
highest-value locations $\{\ell_1, \ldots, \ell_4\}$ and leaves the
lowest-value location $\ell_5$ (Diego Garcia, $v_{\ell_5} = 0.78$) empty.
This is a direct consequence of greedy's myopic value-maximization: it has
no coverage constraint.

Between time steps, sustainment actions are determined by a
\emph{ReplenishmentPolicy} $\pi_R$ that inspects each asset independently:
\begin{equation}
    \pi_R(a) =
    \begin{cases}
        \textsc{Maintain}  & \text{if } r_a < 0.4 \text{ or } d_a < 7, \\
        \textsc{Resupply}  & \text{if } q_a < 2, \\
        \textsc{Hold}      & \text{otherwise,}
    \end{cases}
    \label{eq:policy}
\end{equation}
where $r_a \in [0,1]$ is the readiness rate, $d_a$ is days until next
scheduled maintenance, and $q_a$ is asset quantity.
\textsc{Maintain} restores $r_a \leftarrow \min(1, r_a + 0.20)$;
\textsc{Resupply} increments $q_a \leftarrow q_a + 2$.
The per-action sustainment costs are
$\kappa(\textsc{Reposition}) = 10$,
$\kappa(\textsc{Resupply}) = 5$,
$\kappa(\textsc{Maintain}) = 2$, and
$\kappa(\textsc{Hold}) = 0$.

%------------------------------------------------------------
\subsection{Performance Metrics}
\label{sec:exp1:metrics}

Four metrics are computed at each time step.

\paragraph{Readiness Score ($R$).}
Quantity-weighted mean readiness across all assets:
\begin{equation}
    R = \frac{\sum_{a \in \mathcal{A}} q_a \, r_a}{\sum_{a \in \mathcal{A}} q_a}.
    \label{eq:readiness}
\end{equation}

\paragraph{Coverage Score ($C$).}
Fraction of locations occupied by at least one asset:
\begin{equation}
    C = \frac{|\{\ell \in \mathcal{L} : \exists\, a,\, \pi(a) = \ell\}|}{|\mathcal{L}|}.
    \label{eq:coverage}
\end{equation}

\paragraph{Sustainment Cost ($K$).}
Total action cost incurred by the ReplenishmentPolicy at a given step:
\begin{equation}
    K = \sum_{a \in \mathcal{A}} \kappa\!\left(\pi_R(a)\right).
    \label{eq:cost}
\end{equation}

\paragraph{Posture Efficiency ($E$).}
A composite metric that rewards joint readiness--coverage while penalizing
cost superlinearly:
\begin{equation}
    E = \frac{R \cdot C}{\log(K + 2)}.
    \label{eq:efficiency}
\end{equation}
The logarithmic cost penalty reflects diminishing marginal returns to
sustainment spending.

\paragraph{Scenario-Weighted Readiness (SWR).}
To quantify threat-environment sensitivity, we evaluate greedy postures
against a set of $S = 20$ independently drawn threat scenarios. Each
scenario $s$ assigns a threat level $\tau_\ell^{(s)} \in [0,1]$ to every
location, reducing the effective readiness of any asset stationed there.
The SWR is the probability-weighted expected effective readiness:
\begin{equation}
    \mathrm{SWR} = \sum_{s=1}^{S} \frac{w_s}{\sum_{s'} w_{s'}}
    \cdot \frac{\sum_{a \in \mathcal{A}} q_a \, r_a
    \left(1 - \tau_{\pi(a)}^{(s)}\right)}{\sum_{a \in \mathcal{A}} q_a},
    \label{eq:swr}
\end{equation}
where $w_s = 1$ for all $s$ (uniform weighting).
Two threat distributions are evaluated:
\begin{itemize}
    \item \textbf{Uniform}: $\tau_\ell^{(s)} \sim \mathcal{U}(0.10, 0.30)$
          for all $\ell$---a low-intensity, operationally symmetric threat
          environment.
    \item \textbf{Skewed}: $\tau_\ell^{(s)} = \min\!\left(0.9,\;
          v_\ell \cdot \mathcal{U}(0.50, 1.00)\right)$---threat intensity
          scales with strategic value, modeling an adversary that preferentially
          targets the most valuable basing locations.
\end{itemize}
All scenario draws are fixed by a common seed across simulation runs so
that observed variance reflects initial-state heterogeneity, not
scenario sampling.

%------------------------------------------------------------
\subsection{Results}
\label{sec:exp1:results}

\subsubsection{Per-Step Metrics}

Table~\ref{tab:exp1_metrics} reports the four metrics over all 10 time
steps. Several features merit attention.

\begin{itemize}
    \item \textbf{Readiness declines then stabilizes.} $R$ falls from
    $0.712 \pm 0.057$ at $t = 0$ to a trough near $0.553$ at $t = 5$--$6$,
    then partially recovers to $0.570 \pm 0.041$ at $t = 10$ as the
    ReplenishmentPolicy's \textsc{Maintain} actions restore degraded assets.
    The stabilization reflects an equilibrium between the fixed $\delta = 0.08$
    degradation rate and the $+0.20$ readiness gain from \textsc{Maintain}.

    \item \textbf{Coverage is constant at 0.80.} Greedy's deterministic
    assignment to the four highest-value locations fixes $C = 4/5 = 0.80$
    throughout the simulation. The fifth location remains unoccupied at
    every time step.

    \item \textbf{Sustainment cost exhibits a large initial surge.}
    $K = 13.1 \pm 6.0$ at $t = 0$, reflecting widespread maintenance
    and resupply needs in the initial asset population (low maintenance
    windows, low quantities). Cost drops sharply to $3.4 \pm 1.9$ by
    $t = 2$ as immediate needs are addressed, then rises gradually to
    $6.8 \pm 3.2$ by $t = 10$ as ongoing degradation generates a
    steady maintenance demand.

    \item \textbf{Posture efficiency peaks early then decays.}
    $E$ rises sharply from $0.221$ at $t = 0$ to $0.344$ at $t = 1$
    as the initial cost surge subsides, then declines monotonically
    to $0.222$ by $t = 10$ as readiness erodes and cost accumulates.
    The return to near-$t_0$ efficiency by $t = 10$ indicates that
    greedy reaches a long-run steady state.
\end{itemize}

\begin{table}[htbp]
\centering
\caption{Greedy placement baseline metrics over 10 time steps (95\% CI across 10 random seeds; $\delta = 0.08$; location capacity $= 5$). Values shown as mean~(lower, upper).}
\label{tab:exp1_metrics}
\begin{tabular}{ccccc}
\hline
Step & Readiness & Coverage & Sustainment Cost & Posture Efficiency \\
\hline
         0 & 0.712~(0.671, 0.753) & 0.80 & 13.100~(8.776, 17.424) & 0.221~(0.191, 0.250) \\
         1 & 0.671~(0.638, 0.705) & 0.80 & 4.200~(2.328, 6.072) & 0.344~(0.241, 0.447) \\
         2 & 0.624~(0.599, 0.650) & 0.80 & 3.400~(2.043, 4.757) & 0.337~(0.242, 0.433) \\
         3 & 0.587~(0.562, 0.613) & 0.80 & 4.000~(2.216, 5.784) & 0.304~(0.216, 0.393) \\
         4 & 0.575~(0.553, 0.597) & 0.80 & 5.000~(3.063, 6.937) & 0.252~(0.221, 0.283) \\
         5 & 0.554~(0.526, 0.582) & 0.80 & 5.200~(3.269, 7.131) & 0.239~(0.210, 0.267) \\
         6 & 0.553~(0.524, 0.582) & 0.80 & 5.200~(3.269, 7.131) & 0.238~(0.211, 0.264) \\
         7 & 0.562~(0.533, 0.591) & 0.80 & 5.600~(3.599, 7.601) & 0.236~(0.207, 0.264) \\
         8 & 0.559~(0.528, 0.590) & 0.80 & 5.800~(3.727, 7.873) & 0.233~(0.202, 0.264) \\
         9 & 0.569~(0.535, 0.602) & 0.80 & 6.200~(4.127, 8.273) & 0.231~(0.200, 0.261) \\
        10 & 0.570~(0.541, 0.599) & 0.80 & 6.800~(4.543, 9.057) & 0.222~(0.195, 0.249) \\
\hline
\end{tabular}
\end{table}

\subsubsection{Comparison with Random Placement}

To contextualize greedy's performance, we compare it against a random
placement baseline that assigns assets uniformly at random subject to the
same capacity constraint, followed by the same ReplenishmentPolicy.
Figure~\ref{fig:exp1_main} (panels a--b) shows the trajectories for both
strategies across all 10 seeds.

Readiness is \emph{placement-independent}: both greedy and random converge
to $R = 0.570 \pm 0.041$ by $t = 10$ (panel~a). The ReplenishmentPolicy
applies identically regardless of geographic assignment, so asset readiness
is governed solely by the degradation--repair equilibrium. The two lines
overlap throughout, confirming that value-based greedy placement confers
no readiness advantage over random assignment.

Posture efficiency tells a different story (panel~b). Random placement
achieves $E = 0.278 \pm 0.047$ at $t = 10$, versus $E = 0.222 \pm 0.038$
for greedy---a \textbf{25.1\% gap}. The sole driver is coverage: random
assignment distributes assets across all five locations ($C = 1.00$),
while greedy concentrates assets at four ($C = 0.80$). Given identical
readiness and cost, Equation~\eqref{eq:efficiency} yields a coverage
ratio of $1.00 / 0.80 = 1.25$, fully accounting for the observed
efficiency gap. This result exposes a fundamental limitation of value-only
greedy placement: by ignoring geographic coverage, it sacrifices composite
operational efficiency despite maintaining strong local asset readiness.

\subsubsection{Threat-Environment Sensitivity}

Table~\ref{tab:exp1_sensitivity} and Figure~\ref{fig:exp1_main} (panel~c)
report SWR under uniform and skewed threat distributions. The results
reveal a structural vulnerability of greedy placement.

Under \textbf{uniform threat}, SWR declines from $0.568 \pm 0.046$ at
$t = 0$ to $0.455 \pm 0.032$ at $t = 10$---a 20.2\% reduction from
nominal readiness. This modest degradation reflects the low mean threat
level ($\mathbb{E}[\tau] = 0.20$) under the uniform distribution, which
reduces effective readiness by a corresponding fraction regardless of
where assets are placed.

Under the \textbf{skewed distribution}, SWR collapses to
$0.243 \pm 0.020$ at $t = 0$ and $0.194 \pm 0.015$ at $t = 10$---a
65.9\% reduction from nominal, or a 57.3\% drop relative to the
uniform-threat SWR. This gap is \emph{not a statistical artifact}: it
remains constant at $57.3\%$ across all 10 time steps and all 10 random
seeds (Table~\ref{tab:exp1_sensitivity}, \% Drop column).

The constancy has a structural explanation. Because the threat scenarios
are fixed and greedy \emph{always} assigns assets to the same four
highest-value locations $\{\ell_1, \ldots, \ell_4\}$, the SWR ratio
between conditions is determined entirely by the ratio of expected
survival rates at those locations:
\begin{equation}
    \frac{\mathrm{SWR}_\mathrm{skewed}}{\mathrm{SWR}_\mathrm{uniform}}
    \approx
    \frac{\frac{1}{4}\sum_{i=1}^{4}
          \mathbb{E}\!\left[1 - \tau_{\ell_i}^{(s)}\right]_\mathrm{skewed}}
         {\mathbb{E}\!\left[1 - \tau^{(s)}\right]_\mathrm{uniform}}
    = \frac{0.344}{0.800} \approx 0.43,
    \label{eq:structural_gap}
\end{equation}
yielding a fixed 57\% drop. Under the skewed distribution, an adversary
that targets high-strategic-value bases finds that \emph{all} greedy assets
are concentrated precisely at the locations facing the greatest threat.
Greedy placement, by optimizing for peacetime strategic value, maximally
exposes assets to value-correlated adversarial threats.

\begin{table}[htbp]
\centering
\caption{Greedy scenario-weighted readiness (SWR) under uniform and skewed threat distributions (95\% CI across 10 random seeds; 20 scenarios per condition). $\Delta = \text{SWR}_{\text{uniform}} - \text{SWR}_{\text{skewed}}$; \%~drop $= \Delta / \text{SWR}_{\text{uniform}}$. Values shown as mean~(lower, upper).}
\label{tab:exp1_sensitivity}
\begin{tabular}{ccccc}
\hline
Step & SWR (uniform) & SWR (skewed) & $\Delta$ & \% Drop \\
\hline
         0 & 0.568~(0.536, 0.601) & 0.243~(0.228, 0.258) & 0.325 & 57.3\% \\
         2 & 0.499~(0.478, 0.519) & 0.213~(0.204, 0.222) & 0.286 & 57.3\% \\
         4 & 0.459~(0.441, 0.477) & 0.196~(0.188, 0.203) & 0.264 & 57.4\% \\
         6 & 0.442~(0.419, 0.465) & 0.188~(0.178, 0.199) & 0.254 & 57.4\% \\
         8 & 0.447~(0.422, 0.471) & 0.190~(0.179, 0.202) & 0.256 & 57.4\% \\
        10 & 0.455~(0.432, 0.478) & 0.194~(0.184, 0.205) & 0.261 & 57.3\% \\
\hline
\end{tabular}
\end{table}

\begin{figure}[htbp]
    \centering
    \includegraphics[width=\linewidth]{exp1_main}
    \caption{Greedy baseline characterization over 10 time steps
             ($N_\mathrm{seed} = 10$, $\delta = 0.08$, capacity $= 5$).
             Shaded bands show $\pm 1\sigma$.
             \textbf{(a)}~Readiness score: greedy and random placement
             converge to identical readiness by $t = 5$, confirming
             placement-invariance under the ReplenishmentPolicy.
             \textbf{(b)}~Posture efficiency: random placement consistently
             outperforms greedy by 25\% due to full geographic coverage
             ($C = 1.00$ vs.\ $C = 0.80$).
             \textbf{(c)}~Scenario-weighted readiness (SWR) under uniform
             vs.\ skewed threat: the 57.3\% gap between conditions is
             constant across all time steps, reflecting a structural
             vulnerability of greedy placement to value-correlated
             adversarial threats.}
    \label{fig:exp1_main}
\end{figure}

%------------------------------------------------------------
\subsection{Discussion}
\label{sec:exp1:discussion}

Experiment~1 yields three findings that collectively motivate the
scenario-weighted optimizer evaluated in Experiment~2.

\paragraph{Finding 1: Readiness is placement-invariant under rule-based sustainment.}
The ReplenishmentPolicy equalizes readiness across placement strategies by
step 10. This means greedy's advantage over random is not readiness but
other posture qualities---and in fact greedy underperforms random on the
composite efficiency metric. Any optimizer that claims readiness improvements
over greedy must do so through sustainment policy changes, not placement alone.

\paragraph{Finding 2: Greedy sacrifices geographic coverage for value concentration.}
With capacity $c = 5$ and $|\mathcal{A}| = 20$ assets, greedy leaves one
location persistently uncovered, incurring a permanent 25\% posture
efficiency penalty. This is not a boundary artifact: for any configuration
where $|\mathcal{A}|$ is not a multiple of $|\mathcal{L}| \cdot c$,
greedy will under-cover lower-value locations. A stochastic optimizer that
accounts for coverage through threat scenario diversity can exploit this gap.

\paragraph{Finding 3: Greedy placement creates a structurally fixed threat vulnerability.}
The 57.3\% SWR degradation under skewed threat is not reducible through
better sustainment---it is a geometric consequence of where assets are placed.
The vulnerability is invariant to seed, time step, and readiness level, meaning
it cannot be repaired reactively; it must be anticipated at the placement stage.
This finding directly motivates the key research question in Issue~\#10 (Q1):
an optimizer that is aware of threat distribution heterogeneity at planning
time should be able to trade some value concentration for threat-robustness,
capturing positive EVSS (Expected Value of the Stochastic Solution)
over greedy.

\end{document}
